% P04: replaces the whole body of
% \subsection{Mono- and Bi-Pulley Tendon Routing}\label{sec:disc_pulley}
\subsection{Mono- and Bi-Pulley Tendon Routing}\label{sec:disc_pulley}

The pulley algorithms extend the route search space by allowing the artificial tendon to be reeved over pulleys before insertion, and the naming refers to how many degrees of freedom dominate the routing calculation: the mono-pulley search computes the route on the knee's dominant degree of freedom, and the bi-pulley search computes it on the two dominant degrees of freedom of the biarticular evaluation planes. The purpose of the reversed pulley is a trade of force for range of motion: the tackle halves the insertion-side force while doubling the insertion travel, which can extend the angle range over which a route meets its torque requirement. In the knee-flexor route the pulley acts as a redirect in this configuration rather than a force multiplier, and the two optimized routing configurations differ by more than the pulley alone because each was solved under its own objective: peak predicted knee-axis torque moved from \qty{81.23}{\newton\meter} to \qty{85.74}{\newton\meter}, with off-axis torque held to \qty{0.803}{\newton\meter}. Appendix~\ref{sec:app_routing_extensions} records the optimizer diagnostics and the broader route searches.

A subsequent transmission study re-solved actuator sizes for each tackle configuration at five anatomical muscle sites, making the force-for-travel trade explicit in the sizing. The corresponding torque and margin surfaces are shown in Figures~\ref{fig:PulleySurfaces} and~\ref{fig:PulleyMargins_resolved}. These calculations use the moving-force-point routing, with the adjacent joint fixed for the biarticular evaluation planes; they do not establish simultaneous performance across both joints. Physical verification of the pulley configurations remains future work.
